% Transcription — fonds Alexandre Grothendieck, shelfmark 130, batch 1 (pages 1-20).
% Established from the facsimile archives/batches/130/batch-01.pdf (a mirror of
% https://grothendieck.umontpellier.fr/130.pdf), per the skill
% transcribe-grothendieck. The batch file carries no cover sheet: PDF page N
% is page N of the folder (the Montpellier footer prints « 1 » ... « 20 »),
% and the archivists' pencil numbers bottom-left agree wherever legible
% (« 1 », « 2 », « 3 » ... « 14 », « 16 », « 18 »).
%
% Pass: Opus 5.5 (claude-opus-5-5), 2026-09-23 — first pass, unchecked against the pages by a human.
%
% Dating. The inventory dates the folder « [vers 1971] » (copied verbatim
% into \dating). The folder sits in the inventory group « Descentes »
% (126-128), which is dated « [avant 1970] »: the group's date and the
% folder's own do not agree. Recorded here as an observation only.
%
% Thirteen of the twenty pages are transcribed. Absent: page 1, an otherwise
% blank sheet inscribed in his hand, top right, « Formules de Néron » — a
% cover, whose words become the first \section heading (the same words head
% page 3); pages 2, 8, 15 and 20, blank; pages 17 and 19, two leaves of a
% Bourbaki typescript (n° 383, « -VII-10- » and « -VII-9- », on Dynkin
% diagrams), printed (their text shows through pages 16 and 18), with nothing in his
% hand. Page 7 is the summary sheet of the same typescript (it shows through
% page 6); it is kept only for the five pencil lines in its left margin,
% the typescript itself is not transcribed. He did not paginate these
% leaves.
%
% What the batch is.
%   3-6    « Formules de Néron » (his title). A_K an abelian variety over the
%          function field K of a discrete valuation ring V; a function
%          v(X, a) on pairs (divisor algebraically equivalent to 0, 0-cycle)
%          with disjoint supports, characterised by bilinearity, translation
%          invariance, a normalisation by intersection numbers on the Néron
%          model when X = div(f), and boundedness; its integrality
%          properties ((v), (vi): m v(X,a) ∈ Z, m the order of the
%          component group A_0/A_0^0), whence a pairing
%          A*(K) × A(K) → Z/mZ; the relation (vii) with Néron's local symbol
%          (X, a)_v, so (X, a)_v ∈ (1/m)Z; symmetry; in the global case the
%          sum over v gives the Néron-Tate form. Page 6: when the schematic
%          closure of A_η = Pic^0 of the generic fibre in Pic_{X/S} is a
%          smooth group subscheme, it is a Néron model (X regular).
%   7      Pencil list in the margin of the typescript: Coxeter and Cartan
%          matrices, root systems, reflection groups, Dynkin diagrams.
%   9-14   A different run, no title: X a smooth projective surface over an
%          algebraically closed field, a pencil of hyperplane sections with
%          base locus Z, X~ the blow-up, the fibration X~ → P^1 with
%          generic fibre Y_η; Pic(X)' and Pic(X~)' (classes of degree 0 on
%          the generic fibre); the kernel of Pic(X~)' → Pic^0 of the
%          geometric generic fibre (vertical classes), the image of Pic^0(X)
%          as the fixed part of the Jacobian, NS(X)' → J'(k(η)) with
%          J' = J/J^fixe, surjectivity modulo finite groups via Tsen, and the
%          boxed formula ρ(X~) - 2 = ρ(X) + z - 2 = ρ(J'); then (page 14)
%          the question whether the Néron-Tate form Q on J'(k(η)) satisfies
%          Q(i(ξ), i(η)) = -ξ.η on NS(X)'.
%   16, 18 Back to local symbols: f: X → S projective, X regular, S the
%          spectrum of a discrete valuation ring; cycles C, D with
%          complementary dimension whose intersection number i(C.D) should
%          depend only on the generic fibres, as a local symbol
%          <C_η, D_η>_v independent of the model; the case of k cycles; a
%          list of properties of the local symbol (multilinearity, symmetry,
%          invariance, transposition formula, normalisation by div f,
%          continuity, extension of the base field).
%
% Notation fixed for the batch. \underline{a}, \underline{b} for the 0-cycles
% he underlines; D_a, D_\ell for divisors algebraically, linearly equivalent
% to 0, and Z_0^{*}, Z_\ell^{*} for his 0-cycle groups (the Z is drawn curly
% on page 3 only); \mathcal{A} for the Néron model and the schematic
% closure of pages 3 and 6, which he writes with an A barely distinct from
% that of the abelian variety A (the distinction is editorial and said so
% in a note); \underline{\mathrm{Pic}} for his underlined Pic; \tilde{z}
% for the curly letter shaped like a 3 that labels the exceptional curves
% over the points z ∈ Z (pages 10-11), and \xi_0 for the class of Y_0;
% J^{\mathrm{fixe}} for his « J^fixe ». His « V.A. », « v.d. », « pt »,
% « tt », « nb », « alg. », « équiv. », « hyp. », « cf. » are kept. Square
% brackets in the text are his, except those of the editorial section
% headings of pages 9 and 16.
%
% The dating is the inventory's own for the folder, copied verbatim; the
% title in \foldertitle is the archivists' (not bracketed in the inventory).
\documentclass[11pt,a4paper]{article}
% Le préambule commun à toutes les transcriptions du fonds Grothendieck.
%
% Il tient en une idée : **l'incertitude se compose, elle ne se résout pas.**
% Une transcription qui lisse un mot douteux a détruit la seule chose pour
% laquelle on la faisait — un lecteur doit pouvoir savoir, à chaque endroit,
% ce qui était sur le feuillet et ce qui a été ajouté. Les six macros qui
% suivent sont donc l'appareil critique, et non de la décoration.
%
% Les mêmes commandes sont reconnues par `scripts/render.mjs`, qui produit la
% vue de lecture du site. Une macro ajoutée ici doit l'être aussi là-bas,
% faute de quoi le rendu échouera bruyamment — ce qui est voulu : un rendu
% silencieusement faux est pire qu'un rendu absent.

\ProvidesPackage{grothendieck}[2026/08/08 Transcriptions du fonds Grothendieck]

\usepackage{amsmath,amssymb,amsthm}
\usepackage{mathrsfs}
\usepackage{stmaryrd}
% Les diagrammes commutatifs sont partout dans ce fonds : c'est la langue
% même de la géométrie algébrique de Grothendieck, et une transcription qui
% les rendrait en prose ne transcrirait rien.
\usepackage{tikz-cd}
% Français par défaut — c'est la langue des feuillets — mais l'anglais est
% chargé aussi : la traduction et le résumé sont des documents anglais, et la
% typographie française (espaces avant les deux-points) y serait une faute.
% Ces documents basculent par \selectlanguage{english} en tête de corps.
\usepackage[english,french]{babel}
\usepackage{fontspec}
\usepackage{geometry}
\geometry{a4paper,margin=25mm,marginparwidth=28mm}
\usepackage{marginnote}
\usepackage{xcolor}
% ulem, not soul. `soul`'s \st and \ul are built on a character-by-character
% scan that XeLaTeX's font machinery breaks: the document fails with an
% undefined control sequence rather than degrading. ulem does the same two jobs
% and is Unicode-engine safe. `normalem` stops it hijacking \emph, which the
% transcriptions use for Grothendieck's own emphasis.
\usepackage[normalem]{ulem}
\usepackage{hyperref}
\hypersetup{colorlinks=true,linkcolor=black,urlcolor=black,pdfborder={0 0 0}}

% --- Métadonnées du lot -----------------------------------------------------
% Reprises telles quelles par le script de rendu, qui les affiche en tête de
% la vue de lecture. Elles fixent ce que le fichier prétend couvrir : sans
% elles, une transcription flotte sans point d'attache dans un fonds de seize
% mille feuillets.

\newcommand{\folder}[1]{\gdef\grothFolder{#1}}
\newcommand{\batch}[1]{\gdef\grothBatch{#1}}
\newcommand{\leaves}[2]{\gdef\grothFirst{#1}\gdef\grothLast{#2}}
\newcommand{\foldertitle}[1]{\gdef\grothFoldertitle{#1}}
\gdef\grothFolder{?}\gdef\grothBatch{?}\gdef\grothFirst{?}\gdef\grothLast{?}\gdef\grothFoldertitle{}

% --- Le résumé --------------------------------------------------------------
%
% Il ouvre la lecture modernisée, sous le titre « Résumé », et il est composé
% en petit. Ce n'est pas un ornement : c'est le seul passage du document qui
% ne soit pas des mathématiques, et un lecteur doit pouvoir le reconnaître
% avant de l'avoir lu — puis le sauter s'il connaît déjà le sujet, ou s'y
% arrêter s'il ne le connaît pas. Le filet qui le suit marque la reprise du
% corps.

\newenvironment{resume}
  {\small\setlength{\parskip}{0.45em}}
  {\par\medskip\hrule\medskip}

% --- Mots-clés ---------------------------------------------------------------
%
% En anglais, en clôture du résumé de la lecture modernisée : le vocabulaire
% moderne sous lequel chercher ce que les feuillets construisent. Le manifeste
% du site (`npm run manifest`) les extrait pour étiqueter la cote entière —
% cette ligne est la source unique des tags, il n'y a pas de fichier à côté.
\newcommand{\keywords}[1]{\par\smallskip\noindent
  {\footnotesize\textsc{Keywords} --- \textit{#1}}\par}

% --- L'appareil critique ----------------------------------------------------

% Le numéro de feuillet, en marge. C'est l'ancre qui permet de régler un
% désaccord sur une formule en regardant *un* feuillet plutôt qu'une chemise
% de six cents. Le site s'en sert aussi pour tourner les pages du fac-similé
% quand on fait défiler la transcription.
\newcommand{\leaf}[1]{%
  \marginnote{\footnotesize\color{gray!70}\textsf{#1}}%
  \label{leaf:#1}\ignorespaces}

% Illisible : on ne devine pas. Un mot inventé qui se lit comme les autres est
% le pire résultat possible.
\newcommand{\ill}{\textcolor{red!60!black}{[\,\ldots\,]}}

% Lecture douteuse : le mot est donné, et signalé comme incertain.
\newcommand{\uncertain}[1]{\uline{#1}}

% Ajout de l'éditeur — une lettre manquante, un mot sous-entendu.
\newcommand{\add}[1]{\textcolor{blue!50!black}{[#1]}}

% Biffé par Grothendieck lui-même. Ce qu'il a rayé fait partie du document :
% une rature dit souvent où la pensée a changé de direction.
\newcommand{\struck}[1]{\sout{#1}}

% Note du transcripteur, sur l'état matériel du feuillet.
\newcommand{\note}[1]{\par\smallskip{\small\color{gray!60!black}\textit{[#1]}}\par\smallskip}

% Note marginale de Grothendieck, distincte du corps du texte.
\newcommand{\marginal}[1]{\par\smallskip\noindent\hspace{1em}%
  {\small\color{gray!60!black}#1}\par\smallskip}

% --- Titre ------------------------------------------------------------------

\AtBeginDocument{%
  \begin{center}
    {\large\bfseries Fonds Alexandre Grothendieck --- cote n\textsuperscript{o}~\grothFolder}\\[2pt]
    {\itshape\grothFoldertitle}\\[4pt]
    {\small feuillets \grothFirst--\grothLast\ (lot \grothBatch)}\\[2pt]
    {\footnotesize\color{gray!60!black}Transcription établie d'après le fac-similé numérisé
     par l'Université de Montpellier.\\
     Les lectures incertaines sont soulignées, les illisibles notées [\,\ldots\,],
     les ajouts entre crochets.}
  \end{center}
  \vspace{1em}}

\endinput


\folder{130}
\batch{1}
\pages{1}{20}
\dating{[vers 1971]}
\watermark{Édition de démonstration}
\foldertitle{Hauteurs et symboles locaux : notes manuscrites (s.d.).}

\begin{document}

\section{Formules de Néron}
\note{titre de sa main, en haut à droite du feuillet 1, par ailleurs
vierge, qui ne reçoit pas de numéro de page ; les mêmes mots, soulignés,
ouvrent le feuillet 3}

\page{3}
\emph{Formules de Néron}

$A_K$ V.A. sur $K$ corps des fonctions de $V$, anneau de \uncertain{v.d.}
discrète (corps résiduel fini ?)

On définit une fonction $v(X, \underline{a})$ [à valeurs dans
$\mathbb{Q}$] sur la partie de
\[
  D_a(A_K) \times Z_0^{*}(A)_K
\]
formée des couples $X, \underline{a}$ tels que
$\mathrm{Supp}\, X \cap \mathrm{Supp}\, \underline{a} = \emptyset$,
uniquement déterminée par les propriétés suivantes :
\note{« à valeurs dans $\mathbb{Q}$ » est écrit au-dessus de la ligne,
avec un trait qui le rattache à $v(X, \underline{a})$ ; le $Z$ est ici
une lettre bouclée ; l'exposant, lu $*$, est douteux, et l'indice $K$ est
écrit par-dessus un signe biffé}

(i) Bilinéaire.

(ii) Invariant par translation par les $a \in A(K)$.

(iii) Si $X \sim_{\ell} 0$, \struck{Aus} \struck{et} $X = \mathrm{div}(f)$,
\uncertain{on a}
\[
  v(X, \underline{a}) = \sum_{x_0 \in A_0} i(\mathrm{div}_{g}(\bar{f}), \bar{\underline{a}} ; x_0)
\]
\note{au-dessus de « div », un mot biffé, illisible ; l'indice de
« div », lu $g$, est douteux, et c'est le même signe deux lignes plus bas}

$\mathcal{A}$ est le \emph{modèle de Néron} de $A$,
$A_0 = \mathcal{A} \otimes_V k$, $\bar{f} \to f$ considéré comme fonction
sur $\mathcal{A}$, et $\mathrm{div}_{g} \bar{f} \to$ la composante
verticale de $\mathrm{div}(\bar{f})$, $\bar{\underline{a}} \to$
l'« adhérence » de $\underline{a}$ dans $\mathcal{A}$.
\note{ces trois lignes se ferment sur une accolade ; le modèle de Néron
est écrit d'un $A$ à peine plus grand que celui de la variété abélienne,
et noté ici $\mathcal{A}$}

\marginal{Si $m = 1$ i.e. $A_0$ connexe, alors $v(X, \underline{a}) = 0$
si $X \underset{\mathrm{alg}}{\sim} 0$ et pour \ill{} ; \uncertain{vraiment}
$v(X, \underline{a}) = 0$ $\forall\, X, \underline{a}$}
\note{note oblique dans la marge de gauche, à la hauteur de (iii), reliée
par un trait à « (iii) » ; « connexe » est coupé au bord}

(iv) Pour $X \in D_a(A_K)$ fixé, $0 \notin \mathrm{Supp}\, X$, et
$x \in A(K)$, $x \notin \mathrm{Supp}\, X$ variable, $v(X, (x) - (0))$ est
borné.

De plus, on a les propriétés suivantes :

(iv bis) $v(X, (x) - (a))$ ne dépend (pour $X \in D_a(A_K)$ \struck{fixé}
et $a \in A(K)$, $a \notin \mathrm{Supp}\, X$ fixés) \struck{$x$ variable}
que de la classe de $x$ mod $A(K)^{0} =$ Image inverse de $A_0^{0}$\ill{},
donc ne prend qu'un nb fini de valeurs.
\note{après $A_0^{0}$, un indice coupé au bord du feuillet}

\page{4}
\struck{\ill{}}
\note{en haut à gauche, dans la marge, un griffonnage biffé où se lit
$X$}

(v) Si $X \in D_{\ell}(A_K)$ ou $\underline{a} \in Z_{\ell}^{*}(A_K)$, on a
$v(X, \underline{a}) \in \mathbb{Z}$.

(vi) Si $m =$ degré de $A_0 / A_0^{0}$, alors on a toujours
$m\, v(X, \underline{a}) \in \mathbb{Z}$.

[N.B. On obtient donc [par (v) et (vi)] un accouplement remarquable

\begin{tikzcd}[column sep=small]
  D_a(A_K)/D_{\ell}(A_K) \times Z_0^{*}(A_K)/Z_{\ell}^{*}(A_K) \arrow[r] \arrow[d, no head, "\wr" description] & \mathbb{Z}/m\mathbb{Z} \arrow[d, no head, "\wr" description] \\
  A^{*}(K) & {}_{m}(\mathbb{Q}/\mathbb{Z})
\end{tikzcd}

\note{le $\wr$ de gauche est écrit sous $D_a(A_K)/D_{\ell}(A_K)$ seul ;
le $K$ de $A^{*}(K)$ est surchargé}

\struck{pour \ill{} unicité (?)} pour \ill{}

* $\quad A^{*}(K) \times A(K) \longrightarrow \mathbb{Z}/m\mathbb{Z}$ ]

\marginal{À examiner : cas d'une \uncertain{pol.} \uncertain{principale},
\uncertain{jacobienne} d'une VA et *}
\note{note entourée, oblique, dans la marge de gauche ; l'astérisque
renvoie à la ligne marquée *}

La relation avec le symbole local $(X, \underline{a})_v$ est donnée par

(vii)
\[
  \boxed{(X, \underline{a})_v = v(X, \underline{a}) + \sum_{x_0 \in A_0} i(\bar{X}, \bar{\underline{a}}, x_0)}
\]
\note{un petit signe, illisible, en exposant de $A_0$ sous le $\Sigma$}

Ceci montre que $(X, \underline{a})_v \in \frac{1}{m}\mathbb{Z}$ pour tt
$X, \underline{a}$, \struck{donc} à fortiori il est \emph{rationnel}, et il
est entier si $X \in D_{\ell}(A_K)$ [On le savait déjà, car on a toujours,
si $X = \mathrm{div} f$, $(X, \underline{a})_v = v(f(\underline{a}))$] ou
si $\underline{a} \in Z_{\ell}^{*}(A_K)$. On trouve donc

\marginal{Se \uncertain{généralise} \uncertain{aussitôt} aux variétés
quelconques ----}
\note{dans la marge de gauche, reliée par une accolade aux trois lignes
« Ceci montre que » à « ou si »}

\page{5}
une \struck{autre} [expression de l'] accouplement * en termes de
$(X, \underline{a})_v$. \struck{\ill{}}

\emph{Cor}
\[
  v(X_{\underline{a}}, \underline{b}) = v(X^{-}_{\underline{b}}, \underline{a})
  \qquad (= v(X_{\underline{b}^{-}}, \underline{a}^{-}))
\]

Dans le cas global [où on a une famille de valuations \emph{discrètes}
satisfaisant à la formule du produit] on obtient un accouplement
\struck{\ill{}}
\[
  D_a^{*}(A_K) \times Z_0^{*}(A_K) \longrightarrow \mathbb{R}
\]
\note{l'exposant de $D_a$, lu $*$, est douteux ; le $Z$ est surchargé}
par \struck{produit} produit des symboles locaux
\[
  (X, \underline{a}) = \sum_v (X, \underline{a})_v
\]
qui ne dépend plus que des classes \struck{de} mod $D_{\ell}(A_K)$ resp.
mod $Z_{\ell}^{*}(A_K)$, et définit donc un accouplement
\[
  A^{*}(K) \times A^{*}(K) \longrightarrow \mathbb{R} .
\]
\note{les deux astérisques sont surchargés, le second plus lourdement ;
on attend $A^{*}(K) \times A(K)$, comme à la p.~4}

On obtient donc : cet accouplement est à valeurs dans
$\frac{1}{m}\mathbb{Z} \subset \mathbb{Q}$, où $m$ est un entier $> 0$
convenable.

La forme en question est la \emph{forme de Néron-Tate}.
\struck{Elle \ill{} multipliée}

\page{6}
$X \xrightarrow{f} S$ morphisme \struck{régulier} propre

$X$ régulier

$S$ régulier irréductible, pt. gén. $\eta$

On suppose que $\underline{\mathrm{Pic}}_{X/S}$ existe, et
\struck{est une V.A.}
\note{l'exposant de $\underline{\mathrm{Pic}}$ est biffé}

\struck{Pour} que l'adhérence schématique $\mathcal{A}$ de
$A_\eta = \underline{\mathrm{Pic}}^{0}_{X_\eta/k(\eta)}$, \struck{est} un
\emph{sous-schéma en groupes} \uncertain{plat} de $X$, \struck{en quoi}
\emph{lisse} sur $S$ --- ce qui est automatiquement le cas si
$\underline{\mathrm{Pic}}^{0}_{X/S}$ est lisse sur $S$ le long de la
section unité, p.ex. si $S$ régulier de dim 1 et les
$\underline{\mathrm{Pic}}_{X_s/k(s)}$ sont lisses ($s \in S$).
\note{« schématique $\mathcal{A}$ » est écrit au-dessus de la ligne, sous
une accolade ; l'exposant de $\underline{\mathrm{Pic}}_{X_\eta/k(\eta)}$,
un $0$ surmonté de deux points, est lu $0$}

Alors \struck{Pic} $\mathcal{A}$ est un \emph{modèle de Néron} :

en effet, si $S'/S$ est lisse, alors $X' = X \times_S S'$ est lisse sur
$X$ régulier, donc est régulier, donc \uncertain{tt} \struck{section}
\uncertain{élément} de $\mathrm{Pic}(X'_{\eta'})$ provient d'un
\uncertain{élément} de $\mathrm{Pic}(X')$.

\marginal{Attention s'il n'y a pas de section locale \uncertain{(ét.)} de
$X$ sur $S$}
\note{dans la marge de gauche, à la hauteur de « en effet »}

\page{7}
\note{feuillet dactylographié, non transcrit : sommaire d'une rédaction
Bourbaki, n° 383, chapitre III, « Systèmes de racines (État 3) ». Seules
les cinq lignes au crayon de la marge de gauche sont transcrites ; rien
ne permet d'assurer qu'elles sont de sa main}

\marginal{Matrices de Coxeter ; Matrices de Cartan ; Systèmes de racines ;
groupes engendrés par des réflexions ; diagrammes de Dynkin}
\note{cinq lignes superposées, séparées ici par des points-virgules}

\section{[Groupe de Néron-Severi d'une surface et pinceau de sections hyperplanes]}
\note{titre entre crochets, de l'éditeur : les feuillets 9 à 14 ne portent
pas de titre}

\page{9}
$X$ surface algébrique projective non singulière conn. sur $k$ alg. clos.
\marginal{n'est pas une restriction essentielle ! \struck{\ill{}}
$\dim X \geqslant 2$}
\note{note entourée, en tête du feuillet, reliée par un trait à
« surface »}

$Y_0$ section hyperplane non singulière, dans un « pencil » défini par
$L^{n-2} \subset \mathbb{P}^{n}$ coupant $X$ transversalement en $Z$

$\tilde{X}$ obtenu en faisant éclater $Z$

\begin{tikzcd}
  X & \tilde{X} \arrow[l, "\varphi"'] \arrow[d] & \tilde{Y}_{\bar{t}} \\
  & \mathbb{P}^{1} &
\end{tikzcd}

\note{sous $X$, un trait vertical ondulé, sans pointe ni but ; sous
$\tilde{Y}_{\bar{t}}$, un crochet qui revient vers la gauche, sans
aboutir ; l'indice $\bar{t}$ est douteux}

\begin{tikzcd}
  \mathrm{Pic}(X) \arrow[r, hook] \arrow[dr] & \mathrm{Pic}(\tilde{X}) \arrow[d] \\
  & \mathrm{Pic}(\tilde{Y}_{\bar{\eta}})
\end{tikzcd}

Soit $\mathrm{Pic}(X)'$ ($\mathrm{Pic}(\tilde{X})'$) le sous-groupe des
\struck{diviseurs} classes induisant sur $\tilde{Y}_{\bar{\eta}}$ une classe
de \emph{degré} 0. On a donc

\begin{tikzcd}
  \mathrm{Pic}(X^{*})' \arrow[r, hook] \arrow[dr] & \mathrm{Pic}(\tilde{X})' \arrow[d] \\
  & \mathrm{Pic}^{0}(\tilde{X}_{\bar{\eta}}) = J^{0}(\tilde{X}_{\bar{\eta}})
\end{tikzcd}

\note{l'exposant de $X$, lu $*$, est douteux et absent de la définition
qui précède ; en bas, il écrit $\tilde{X}_{\bar{\eta}}$, non
$\tilde{Y}_{\bar{\eta}}$, et de même aux pp.~10-11}

\page{10}
\note{croquis : une droite de base portant trois points, notés
\uncertain{$t_1$}, \uncertain{$y_0$} et $\bar{\eta}$ ; au-dessus de chacun
une fibre verticale — celle de $t_1$ se recoupe elle-même, celle de
$\bar{\eta}$ est hachurée —, et une courbe transverse, marquée
$\tilde{z}$, qui les coupe toutes les trois}

À déterminer a) les noyaux de
\[
  \begin{array}{l}
    \mathrm{Pic}(X^{*})' \longrightarrow J^{0}(\tilde{X}_{\bar{\eta}}) \\
    \mathrm{Pic}(\tilde{X})' \longrightarrow J^{0}(\tilde{X}_{\bar{\eta}})
  \end{array}
\]
\note{les deux flèches aboutissent au même but ; le $J$ est écrit
par-dessus « Pic »}

b) leur image

\emph{N.B.} $\mathrm{Pic}(\tilde{X}_{\eta}) \to \mathrm{Pic}(\tilde{X}_{\bar{\eta}})$
est injectif, donc le noyau de
$\mathrm{Pic}(\tilde{X})^{*} \to \mathrm{Pic}(\tilde{X}_{\bar{\eta}})$
\struck{\ill{}} (égal au noyau de
$\mathrm{Pic}(\tilde{X})' \to J^{0}(\tilde{X}_{\bar{\eta}})$) est aussi
celui de $\mathrm{Pic}(\tilde{X}) \to \mathrm{Pic}(\tilde{X}_{\eta})$, i.e.
est formé des classes des \emph{diviseurs verticaux}. \struck{Parmi
ceux-ci,} Notons que les diviseurs verticaux $\tilde{X}_y$, pour
$y \in \mathbb{P}^{1}(k)$, \struck{/ $\tilde{X}_y$ non sing\ill{}} sont
linéairement équivalents entre eux, donc à
\[
  \tilde{X}_{y_0} = \tilde{Y}_0 \simeq Y_0 .
\]
\struck{Donc} Si \struck{les} tous les $X_y$ sont \struck{irréductibles}
\struck{\ill{}} \emph{intègres}, alors le noyau de
$\mathrm{Pic}(\tilde{X})' \to J^{0}(\tilde{X}_{\bar{\eta}})$ se réduit aux
\struck{classes des} multiples de \struck{cl($\tilde{X}_y$)} la classe
\struck{\ill{}} de $\tilde{Y}_0$, \uncertain{notée} $\tilde{\xi}_0$
\struck{Notons que}, qui est
\note{« notée $\tilde{\xi}_0$ » est écrit au-dessus de « Notons que »
biffé}

\page{11}
i.e. à
\[
  \varphi^{*}(\xi_0) - \sum_{z \in Z} \tilde{z} ,
\]
où $\xi_0$ est la classe de $Y_0$ dans $X$. Par suite, \struck{si}
\uncertain{comme} $Z \neq \emptyset$, on a
$n \in \mathbb{Z}$, $n \neq 0 \Rightarrow n\tilde{\xi}_0 \notin \mathrm{Pic}(X)'$.
Donc dans le cas envisagé,
\[
  \mathrm{Pic}(X)' \longrightarrow \mathrm{Pic}(\tilde{X}_{\bar{\eta}}) = \mathrm{Pic}(Y_{\bar{\eta}})
\]
est \emph{injectif}.
\note{« comme » est écrit au-dessus de « si » biffé ; le prime de
$\mathrm{Pic}(X)'$ est appuyé}

L'image de $\mathrm{Pic}^{0}(X) \simeq \mathrm{Pic}^{0}(\tilde{X})$ dans
\struck{$\mathrm{Pic}(\tilde{X}_{\bar{\eta}})$}
$\mathrm{Pic}^{0}(Y_{\bar{\eta}}) = J(Y_{\bar{\eta}})(k)$ est la « partie
fixe » $J(Y_{\bar{\eta}})^{\mathrm{fixe}}(k)$ et de façon précise on a des
isom.

\begin{tikzcd}
  \mathrm{Pic}^{0}(X) \arrow[r, "\sim"] & \mathrm{Pic}^{0}(\tilde{X}) \arrow[r, "\sim"] & J(Y_{\bar{\eta}})^{\mathrm{fixe}}(k) \arrow[d, no head, "\wr" description] \\
  & & J(Y_{\bar{\eta}})^{\mathrm{fixe}}(k(\eta))
\end{tikzcd}

\note{le premier $J$ est écrit par-dessus « Pic » ; après « $(k$ », une
surcharge biffée}

[\emph{N.B.} \struck{\ill{}} $k(\eta)$ est une extension \emph{pure} de
\uncertain{$k(\eta)$}, donc pour tte V.A. $A$ sur $k$,
$A(k) \simeq A(k(\eta))$]. On obtient donc
\note{la lecture « de $k(\eta)$ » est douteuse ; ce qui suit ferait
attendre « de $k$ »}

\begin{tikzcd}
  NS(X)' \arrow[r, hook] & NS(\tilde{X})' \arrow[d] \\
  & J'(k(\eta))
\end{tikzcd}

où $J' = J / J^{\mathrm{fixe}}_{k(\eta)}$ est « la partie mobile » de
$J = J(Y_\eta / k(\eta))$. Notons que le noyau de

\page{12}
$NS(\tilde{X})' \to J'(k(\eta))$ est égal au groupe engendré par l'image
de $\tilde{\xi}_0$, car si $\alpha$ est dans le noyau, alors $\alpha$
provient de $\alpha_0 \in \mathrm{Pic}(\tilde{X})'$, et comme par hyp.
l'image de $\alpha_0$ dans $J(Y_\eta)$ est dans
$J(Y_\eta)^{\mathrm{fixe}} \simeq \mathrm{Pic}^{0}(\tilde{X})$, il
s'ensuit que quitte à modifier $\alpha_0$ par un élément de
$\mathrm{Pic}^{0}(\tilde{X})$, on aura une image nulle --- Donc sous
réserve que tous les $Y_y$ soient intègres, le noyau de
$NS(\tilde{X})' \to J'(k(\eta))$ est de rang 1, et engendré par l'image
de $\tilde{\xi}_0$.
\note{le $J$ de « dans $J(Y_\eta)$ » est surchargé}

Enfin, je dis que \struck{l'image de} $NS(\tilde{X})' \to J'(k(\eta))$ est
\emph{surjectif} mod. groupe fini. En effet, \struck{un élément de}
\struck{$J'(k(\eta))$} comme $H^{1}(K, J'')$ est un groupe de torsion
[où $K = k(\eta) = k(t)$, $J'' = J^{\mathrm{fixe}}$] --- on voit que mod.
groupes finis, \struck{$NS(\tilde{X})' \to J'(k(\eta))$ est}
$J(k(\eta)) \to J'(k(\eta))$ est surjectif), donc il suffit de prouver que
$\mathrm{Pic}(\tilde{X})' \to J(k(\eta))$ est surjectif \struck{mod}
\note{« mod. groupe fini » est écrit au-dessus de la ligne, avec un trait
qui le rattache à « surjectif »}

\page{13}
\struck{\ill{} mod groupe fini}. Or un élément de $J(k(\eta))$
\struck{qui} provient \struck{\ill{}} \struck{\ill{} de $t$} d'une classe
de diviseurs de degré 0 sur $Y_\eta$ [\emph{N.B.} le groupe de Brauer de
$k(\eta) = k(t)$ est nul, par Tsen] et on gagne.

\[
  \boxed{\rho(\tilde{X}) - 2 =}\ \boxed{\rho(X) + z - 2 = \rho(J')}
\]
\note{deux cadres accolés ; les deux « 2 » sont repassés, le second
lourdement, sans doute sur un autre chiffre}

$\rho(X)$ nb de Picard de $X$

$z =$ nb des \uncertain{sections singulières} du pinceau

$\rho(J')$ \struck{=} rang du groupe des pts rat. sur $k(\eta)$, de
$J' = J(Y_\eta) / J(Y_\eta)^{\mathrm{fixe}}$
\note{les trois définitions sont réunies par une accolade ; « sections
singulières » est une lecture douteuse}

Condition de validité de la formule : toutes les sections du pinceau sont
\struck{intègres} \emph{irréductibles} (intègre n'est pas absolument
\uncertain{néc.}, i.e. on \uncertain{admet} des sections multiples ...).

\page{14}
Supposons $X$ une \emph{surface}.

Notons maintenant que $J'$ a une polarisation canonique, héritée de celle
de $J$, \struck{\ill{}} \struck{[\uncertain{En effet, le dual de} $J$
\ill{}] propre dual\ill{}, $\hat{J}'$ le dual de $J'$ s'envoie
canoniquement dans $J$, donc hérite une polarisation}, d'où une
\uncertain{polarisation} \uncertain{Donc on conclut} [par Néron-Tate] une
forme quadratique canonique \struck{sur} définie positive $\mathbb{Q}$
sur $J'(k(\eta))$.
\note{le passage entre crochets est encadré et barré de quatre grands
traits obliques ; en tête du cadre, $J$ suivi d'une flèche. « par
Néron-Tate » est écrit sous la ligne, avec une accolade ; « $\mathbb{Q}$ »
est ajouté au-dessus, avec un signe d'insertion, entre « positive » et
« sur »}

D'autre part, il y a \struck{aussi} une forme quadratique définie
\emph{négative} sur $NS(X)'$, savoir la forme intersection.

$i$ \struck{\ill{}} : $NS(X)' \longrightarrow J'(k(\eta))$. À prouver la
formule
\[
  \boxed{Q(i(\xi), i(\eta)) = -\, \xi \cdot \eta} \qquad (\xi, \eta \in NS(X)')
\]
\marginal{?}
valable si les $Y_t$ sont tous \emph{intègres} ?

\section{[Symboles locaux et nombres d'intersection]}
\note{titre entre crochets, de l'éditeur : les feuillets 16 et 18 ne
portent pas de titre}

\page{16}
1) $\quad f : X \to S$ projectif, $X$ régulier, $S$ spectre d'anneau de
\uncertain{v.d.} discrète, pt générique $\eta$, pt spécial $s$

$X_\eta$ est lisse

$X_s$ est géométriquement intègre \struck{\ill{}}

\struck{Soit} $\dim X_\eta = n \qquad \dim X = n + 1$

$C$, $D$ deux cycles sur $X$, de dimensions $c$ et $d$ avec
$c + d = n + 1$. On suppose

a) \struck{\ill{}} $C_\eta$ et $D_\eta$ sont alg. équiv. à 0

\struck{\ill{} et $\mathrm{Supp}\, C_\eta \cap \mathrm{Supp}\, D_\eta = \emptyset$}

b) $C$ et $D$ n'ont \uncertain{que des} composantes horizontales, i.e.
$C = \overline{C_\eta}$, $D = \overline{D_\eta}$

c) $C \cdot D$ est défini
[$\Longrightarrow \mathrm{Supp}\, C_\eta \cap \mathrm{Supp}\, D_\eta = \emptyset$]

Considérons $i(C \cdot D)$, il ne dépend que de \struck{\ill{}} $C_\eta$ et
$D_\eta$, cycles de $X_\eta$ de codimensions $c, d$ tels que
$c + d = n + 1$, de supports \struck{\ill{}} disjoints. On veut donc
donner une description de $i(C \cdot D)$ en termes de $C_\eta, D_\eta$,
i.e. comme un
\[
  \langle C_\eta, D_\eta \rangle_v
\]
\emph{indépendant du modèle $X$} [et de leur existence], (à valeurs dans
$\mathbb{Q}$) \marginal{à priori}
\note{« algébriquement » est écrit au-dessus de « codimensions » ; « de
supports » au-dessus d'un mot biffé ; « et de leur existence » au-dessus
de la ligne, sous une accolade ; « à priori » dans une bulle reliée à
« à valeurs dans $\mathbb{Q}$ »}

\page{18}
Le cas gl. peut se formuler pour un système de cycles
$Z_1, \dots, Z_k$ de codim $c_1, \dots, c_k$
\[
  \left\{
  \begin{array}{l}
    c_1 + \dots + c_k = n + 1 \\
    \text{[deux parmi les termes]}\ Z_{i\eta}\ \text{alg. équiv. à } 0 \\
    Z_1 \cdot \ldots \cdot Z_k \ \text{défini}
  \end{array}
  \right.
\]
\note{au-dessus de $Z_{i\eta}$, un court mot entouré, illisible ;
« deux parmi les termes » est entouré aussi}

On \uncertain{aurait ainsi}
\[
  i(Z_1 \cdot \ldots \cdot Z_k) = (Z_{1\eta}, \dots, Z_{k\eta})_v
\]

\emph{Propriétés du symbole local}

(i) Bilinéarité (partielles ....), symétrie

(ii) Invariance par isomorphismes de $X_\eta$

(ii bis) Pour $k = 2$, formule de transposition pour un morphisme
$X_\eta \to Y_\eta$

(iii) \struck{Bilinéarité} Normalisation : si
$Z_{k\eta} = Z'_{k\eta} \cdot \mathrm{div} f_\eta$, \struck{Ou} où
$Z_{1\eta} \cdot \ldots \cdot Z_{k\eta} \cdot Z'_{k\eta}$ est défini (i.e.
l'intersection des $\mathrm{Supp}\, Z_{1\eta}$ ... $\mathrm{Supp}\, Z_{k-1,\eta}$,
$\mathrm{Supp}\, Z'_{k\eta}$ est discrète) on a
\struck{\ill{}} \ill{}
\[
  (Z_{1\eta}, \dots, Z_{k\eta}) = v(f_\eta(Z_{1\eta} \cdot \ldots \cdot Z_{k-1,\eta} \cdot Z'_{k\eta}))
\]
\note{devant la parenthèse de gauche, deux lettres, peut-être « i.e. »
ou l'indice $v$ du symbole, non lues ; le dernier indice est coupé au
bord du feuillet}

(iv) Principe de continuité ....

\marginal{(v) Extension du corps de base,}
\note{écrit de bas en haut le long de la marge de gauche}

\end{document}
